% Generated draft: requires mathematical and source review.
\documentclass[11pt]{article}
\usepackage[margin=25mm]{geometry}
\usepackage{fontspec}
\setmainfont{Latin Modern Roman}
\usepackage{amsmath,amssymb,mathtools,unicode-math}
\setmathfont{Latin Modern Math}
\usepackage{enumitem,needspace,xurl,hyperref}
\setlength{\parindent}{0pt}
\setlength{\parskip}{5pt}
\setlength{\emergencystretch}{3em}
\newcommand{\N}{\mathbb{N}}
\newcommand{\Z}{\mathbb{Z}}
\newcommand{\Q}{\mathbb{Q}}
\newcommand{\R}{\mathbb{R}}
\newcommand{\C}{\mathbb{C}}
\newcommand{\F}{\mathbb{F}}
\newcommand{\A}{\mathbb{A}}
\newcommand{\PP}{\mathbb{P}}
\newcommand{\E}{\mathbb{E}}
\newcommand{\eps}{\varepsilon}
\newcommand{\dd}{\,\mathrm{d}}
\begin{document}
\section*{UnsolvedMath ID: 2301006}
\textbf{Source problem code:} catalogue entry
\subsection{Definitions and mathematical statement}
Let $f$ be a nonconstant entire function. Its Nevanlinna characteristic is $T(r,f)$, and its order is $\rho(f)=\limsup_{r\to\infty}\log T(r,f)/\log r$. For a finite value $a\in\C$, let $N(r,a;f)$ be the integrated counting function of zeros of $f-a$, counted with multiplicity. Its deficiency is
\[
 \delta(a,f)=1-\limsup_{r\to\infty}\frac{N(r,a;f)}{T(r,f)}\in[0,1],
\]
and $a$ is deficient if $\delta(a,f)>0$. Write $D(f)=\{a\in\C:\delta(a,f)>0\}$ for the set of \emph{finite} deficient values. The value $\infty$ is automatically deficient for an entire function, so whether it is included in the source's sets must be checked.

\textbf{Exact deficient-value problem (N. U. Arakelyan).} The source reports that, given $\rho>\tfrac12$ and a countable set $E$, Arakelyan constructed an entire function of order exactly $\rho$ for which every point of $E$ is deficient. It asks whether there can be no additional deficient values. With $E\subset\C$ and finite deficient values intended, the quantified question is
\[
 \forall\rho>\tfrac12\ \forall E\subset\C\text{ countable}\quad
 \exists f\text{ entire}:\rho(f)=\rho\text{ and }D(f)=E.
\]
No prescribed numerical deficiencies are requested.

\textbf{Summability conjecture.} For an entire function of finite order with deficient values $a_n$, the source conjectures
\[
 \sum_n\bigl(\log[1/\delta(a_n,f)]\bigr)^{-1}<+\infty.
\]
The natural summability interpretation lists distinct finite deficient values without repetition. The displayed summands are positive and finite only for $0<\delta(a_n,f)<1$; at $\delta=1$ the denominator is zero. Such deficiencies are not excluded by the source, so their treatment is unresolved. A modified logarithm, deletion of exceptional terms, or an extra hypothesis would change the literal formula and is not imposed here.

The supplied summary labels both parts OPEN-TRIAGE, says they were open in Hayman's 2018 edition, relates the first to the Drasin--Weitsman program and Problem 1.26, and mentions the partial/related bound $\sum\delta^{1/3}<\infty$ attributed to Weitsman. It reports no recent resolution found. These are supplied status and progress claims, not independently verified results.
\subsection{Short English statement}
Can any countable set be exactly the finite deficient-value set of an entire function of prescribed order $\rho>1/2$? The second conjecture asks for logarithmic summability of deficiencies of finite-order entire functions, but its formula requires clarification at deficiency $1$ and about whether infinity is included.
\subsection{Sources}
\begin{itemize}
\item Local UnsolvedMath registry, permanent ID 2301006, source code catalogue entry, ``Research Problems in Function Theory --- Problem 1.6''; attribution: N. U. Arakelyan. Supplied status open, classification OPEN-TRIAGE, updated 2026-08-06.
\item Hayman, \emph{Research Problems in Function Theory} (2018), Problem 1.6, \url{https://arxiv.org/abs/1809.07200}. The record reports source-TeX extraction and access on 2026-07-29; the 2018 update was retained for status triage. Current status and rights require review; the record describes publicly available arXiv source TeX.
\item The supplied statement attributes the known existence result to Arakelyan. The research summary refers to the Drasin--Weitsman program, Problem 1.26 of the same list, and a Weitsman summability result, without complete citations or hypotheses.
\end{itemize}
\end{document}
